\hypertarget{link_to_switch}{}{\Large \bf switch*}

{\it Set the output state to one among $n$ input states, based on the value of a controlling state.

Separate blocks exist for $n=2,3,4$ and $5$}.

\begin{center}
\includegraphics{codegen/switch.pdf}
\end{center}

Syntax : \hspace*{10ex}\parbox[t]{60mm}{\sf 
\& switch2 \\
name of variable $x_i(1)$\\
name of variable $x_i(2)$\\
name of variable $x_k$\\
name of variable $x_j$
}\parbox[t]{60mm}{\sf 
\& switch3 \\
name of variable $x_i(1)$\\
name of variable $x_i(2)$\\
name of variable $x_i(3)$\\
name of variable $x_k$\\
name of variable $x_j$
}

~

\hspace*{18ex}\parbox[t]{60mm}{\sf 
\& switch4 \\
name of variable $x_i(1)$\\
name of variable $x_i(2)$\\
name of variable $x_i(3)$\\
name of variable $x_i(4)$\\
name of variable $x_k$\\
name of variable $x_j$
}
\parbox[t]{60mm}{\sf 
\& switch5 \\
name of variable $x_i(1)$\\
name of variable $x_i(2)$\\
name of variable $x_i(3)$\\
name of variable $x_i(4)$\\
name of variable $x_i(5)$\\
name of variable $x_k$\\
name of variable $x_j$
}


~ \hrule ~

Internal states : none

Discrete variables : $z \in \{1, 2, \dots, n \}$

Equations :
\begin{equation*}
0 = x_j - x_i(z)
\end{equation*}

Discrete transitions : 
\[ z \gets \max \left(1, \min (n, nint(x_k) ) \right) \]
where the $\mbox{nint}$ function returns the nearest integer.

\vspace*{2ex}
Initialization of discrete variables : same code as for the discrete transitions
